\section{Introduction}
\label{sec:intro}

Firebase changed how developers build applications. Before Firebase, building
a real-time collaborative app required provisioning servers, configuring
databases, implementing WebSocket synchronization, building authentication
flows, and managing deployment infrastructure. Firebase collapsed all of
this into a single SDK: import the library, initialize with a project ID,
and start writing application logic. Firestore for data, Firebase Auth for
identity, Cloud Functions for server-side logic, Firebase Hosting for
deployment. The developer experience is excellent.

The developer experience is excellent because the developer has ceded
control over every layer of the stack to Google. Firestore data is stored
on Google's servers, encrypted with Google's keys, accessible through
Google's APIs. Firebase Auth identities are Google-managed---the developer
cannot export the authentication state, the password hashes, or the session
tokens. Cloud Functions execute in Google's runtime, with Google-controlled
cold start latency, memory limits, and pricing. Firebase Hosting serves
static assets from Google's CDN with Google-controlled caching, routing,
and availability.

This is the bargain: zero infrastructure complexity in exchange for
complete infrastructure dependency. The cost becomes apparent when:

\begin{enumerate}[label=(\roman*)]
    \item Google changes pricing (as they did with Firestore reads in 2023).
    \item Google deprecates a feature (as they did with Firebase Dynamic Links in 2025).
    \item A regulation requires data residency that Google's infrastructure
    cannot guarantee.
    \item A user requests data deletion and the developer discovers they
    cannot verify that Google actually deleted the data from all replicas.
    \item The application grows beyond Firebase's pricing tiers and
    migration to a cheaper provider requires rewriting the entire data layer.
\end{enumerate}

Supabase, Convex, Neon, and PlanetScale address some of these concerns
(open-source core, PostgreSQL compatibility, edge deployment) but retain
the fundamental architecture: the provider holds the data, the provider
manages the keys, the provider controls the runtime. Switching from
Firebase to Supabase is a lateral move, not an exit.

This paper presents the actual exit: migrating from cloud-dependent
architectures to local-first cryptographic applications where the user
owns the data, holds the encryption keys, and can switch providers
(or run without any provider at all) without losing access to their
application or its state.

Section~\ref{sec:right} acknowledges what Firebase gets right.
Section~\ref{sec:wrong} identifies what it gets wrong.
Section~\ref{sec:mapping} defines the component mapping.
Section~\ref{sec:migration} provides the migration guide.
Section~\ref{sec:performance} benchmarks performance.
Section~\ref{sec:offline} analyzes offline-first advantages.
Section~\ref{sec:cases} presents case studies.
Section~\ref{sec:related} discusses related work.
Section~\ref{sec:conclusion} concludes.

